\documentclass[10pt,a4paper]{article}
\usepackage[margin=24mm]{geometry}
\usepackage{amsmath,amssymb,amsthm,mathtools,enumitem}
\usepackage{hyperref}
\newtheorem{theo}{Theorem}[section]
\newtheorem{lemm}[theo]{Lemma}
\theoremstyle{definition}
\newtheorem{defi}[theo]{Definition}
\allowdisplaybreaks
\setlength{\emergencystretch}{3em}
\DeclareMathOperator{\Homeo}{Homeo}
\DeclareMathOperator{\Clopen}{Clopen}
\DeclareMathOperator{\Aut}{Aut}
\DeclareMathOperator{\id}{id}


%\newcounter{dummy}
%\makeatletter
%\newcommand\myitem[1][]{\item[#1]\refstepcounter{dummy}
%\newcommand*\@currentlabel{[1]}}
%\makeatother

\newcommand{\actson}{\curvearrowright}


\newcommand{\N}{\mathbf{N}}
\newcommand{\Z}{\mathbf{Z}}

\newcommand{\mcA}{\mathcal{A}}
\newcommand{\mcB}{\mathcal{B}}
\newcommand{\mcC}{\mathcal{C}}




\newcommand*{\mk}{\mkern -1mu}
\newcommand*{\Mk}{\mkern -2mu}
\newcommand*{\mK}{\mkern 1mu}
\newcommand*{\MK}{\mkern 2mu}

\hypersetup{urlcolor=purple, linkcolor=blue, citecolor=red}

\newcommand*{\relabel}{\renewcommand{\labelenumi}{(\theenumi)}}
\newcommand*{\romanenumi}{\renewcommand*{\theenumi}{\roman{enumi}}\relabel}
\newcommand*{\Romanenumi}{\renewcommand*{\theenumi}{\Roman{enumi}}\relabel}
\newcommand*{\alphenumi}{\renewcommand*{\theenumi}{\alph{enumi}}\relabel}
\newcommand*{\Alphenumi}{\renewcommand*{\theenumi}{\Alph{enumi}}\relabel}
\let\oldtilde\tilde
\renewcommand*{\tilde}[1]{\mathchoice{\widetilde{#1}}{\widetilde{#1}}{\oldtilde{#1}}{\oldtilde{#1}}}
\let\oldexists\exists
\renewcommand*{\exists}{\mathrel{\oldexists}}
\let\oldforall\forall
\renewcommand*{\forall}{\mathrel{\oldforall}}

\newcommand{\llbracket}{\mathopen{[\mkern -2.7mu[}}
\newcommand{\rrbracket}{\mathclose{]\mkern -2.7mu]}}
%%%%%%%%%%%%%%%%%%%%%%%%%%%%%%%%%%%%%%%%%%%%%%%%%%%%%%
\newcounter{setofstep}
\newcounter{step}[setofstep]
%\crefname{step}{Step}{Steps}
\newenvironment{step}[1][]{\refstepcounter{step}\begin{proof}[{Step~\thestep\if#1\empty\else: #1\fi}]}{\let\qed\relax\end{proof}}
\newcommand{\newstepset}{\stepcounter{setofstep}}

\renewcommand*{\to}{\mathchoice{\longrightarrow}{\rightarrow}{\rightarrow}{\rightarrow}}

\newcommand*{\To}{\mathchoice{\Longrightarrow}{\Rightarrow}{\Rightarrow}{\Rightarrow}}

\let\oldLeftrightarrow\Leftrightarrow
\renewcommand*{\Leftrightarrow}{\mathchoice{\Longleftrightarrow}{\oldLeftrightarrow}{\oldLeftrightarrow}{\oldLeftrightarrow}}
%%%%%%%%%%%%%%%%%%%%%%%%%%%%%%%%%%%%%%%%%%%%%%%%%%%%%%%%%%%%%%%%%%%%%%%%%%%%%%%%%%%%%%
\usepackage{cleveref}
\crefname{originalsection}{original section}{original sections}
\Crefname{originalsection}{Original section}{Original sections}
\makeatletter
\newlabel{s:Krieger}{{4}{source}{}{sourcepointer-1}{}}
\newlabel{s:Krieger@cref}{{[originalsection][0][]4}{[1][0][]source}{}{}{}}
\makeatother
\title{Realizing a prescribed inclusion of ample groups as a small extension on a Cantor set (Lemma 6.2)}
\author{Julien Melleray}
\date{doi:10.5802/ahl.276; CC BY 4.0}
\begin{document}
\maketitle
\setcounter{section}{1}
\noindent\textit{Editorial note.} One independent Lemma 6.2 construction, no separate counting of background lemmas or construction steps.\par\medskip
\noindent\textit{Editorial note.} Original native mathematics retained; official bibliography reconstructed because external .bib is absent.\par\medskip
\noindent\textit{Editorial scope.} Only Lemma 6.2 counted. Standard wrapper and faithful bibliography reconstruction; no mathematical alteration.\par\medskip
\noindent\textit{Allowed standard prerequisites (editorial).}\begin{itemize}
\item Classical Krieger finite systems and conjugacy theorem. Wolfgang Krieger [Kri80], Math. Ann. 252 (1980), no. 2, 87–95, Lemma 2.1 and Theorem 3.5; DOI 10.1007/BF01420115; native lines 161–184 and 622. Exhaustive systems initialize the construction; equality of clopen orbit relations yields conjugacy in Step 3. The stronger current Theorem 4.1 is not invoked.
\item Stone duality for zero-dimensional compact metrizable spaces. Native lines 567–569 and 591. Nested clopen-algebra maps produce the homeomorphism Y to K; compatible finite permutation groups act on the clopen algebra.
\item Glasner–Weiss comparison and compact invariant-measure tools. Source Lemma 2.2; Glasner–Weiss [GW95], Int. J. Math. 6 (1995), no. 4, 559–579; Dini theorem in source Lemma 3.2. Earlier comparison tool retained as stated; all source-local malleability/induced-group proofs retained, including full Lemma 3.3 argument. No external standard proof invented.
\end{itemize}\medskip
\ifdefined\hypertarget\hypertarget{sourcepointer-1}{}\fi
\noindent\textit{Source reference.} Original section 4: local-coloring-linear-logic/\allowbreak{}AHL276/\allowbreak{}sources/\allowbreak{}author-published.tex, source lines 351--359. Forward historical/contextual pointer to a strengthening. Selected construction uses the classical Krieger theorem, not this strengthening.\par\medskip
\paragraph{Source and scope.} Julien Melleray, \emph{A dynamical proof of Matui's absorption theorem}, Annales Henri Lebesgue 9 (2026), 815--837, \href{https://doi.org/10.5802/ahl.276}{doi:10.5802/ahl.276}. Published version, CC BY 4.0.
Only Lemma 6.2 is selected as one complete construction. Sections 2--3 and Definition 5.1 provide uncounted context and author proofs; Definition 6.1 and all four steps of Lemma 6.2 are retained. Original numbering is preserved.
The classical Krieger theorem is an external prerequisite. The strengthened Theorem 4.1, the absorption theorem 6.3 and the appendix correcting earlier work are not selected and are not used as substitutes for this construction.
The native mathematical text is unchanged. This standard article wrapper and bibliography transcription are editorial presentation changes; no mathematical certification or proof repair is asserted.

\section{Background and vocabulary}

We recall some notions and terminology. The interested reader can find a much more detailed exposition in~\cite{Melleray2023}, with proofs for some statements we only mention here in passing.

Assume that $X$ is a compact, metrizable, $0$-dimensional space. Given a subgroup $\Gamma \le \Homeo(X)$, we denote by $M(\Gamma)$ the set of all Borel probability measures on $X$ which are $\Gamma$-invariant. The \emph{full group} $F(\Gamma)$ generated by $\Gamma$ is the set of all $g \in \Homeo(X)$ such that there exists a clopen partition $X= \bigsqcup_{i=1}^n U_i$ and $\gamma_1,\dots,\gamma_n \in \Gamma$ such that $g(x)=\gamma_i(x)$ for all $x \in U_i$. We say that $\Gamma$ is a full group if $\Gamma=F(\Gamma)$. Note that we always have $M(\Gamma)=M(F(\Gamma))$.

For any countable subgroup $\Gamma \le \Homeo(X)$, we denote by $R_\Gamma$ the associated equivalence relation on $X$, and let $[R_\Gamma]$ denote the subgroup of all $g \in \Homeo(X)$ which map each $\Gamma$-orbit to itself. Then $[R_\Gamma]$ is a full group and $M([R_\Gamma])=M(\Gamma)$.

An \emph{ample group over $X$} is a countable, locally finite full group $\Gamma \le \Homeo(X)$ with the property that for all $\gamma \in \Gamma$ the set $\{ x \in X : \gamma(x)=x \}$ is clopen in $X$.

Given a subgroup $\Gamma \le \Homeo(X)$ and a Boolean subalgebra $\mcA$ of $\Clopen(X)$, we say that $(\mcA,\Gamma)$ is a \emph{unit system}\footnote{There is (again) an imprecision in~\cite{Melleray2023}: the definition given there is different from this one for unit systems which are not finite. This is not an issue since those are not used anywhere in~\cite{Melleray2023}.} if:
\begin{itemize}
\item For every $A \in \mcA$ and every $\gamma \in \Gamma$, $\gamma(A) \in \mcA$, giving us an evaluation map $e_\mcA \colon \Gamma \to \Aut(\mcA)$.
\item The morphism $e_\mcA$ is injective (equivalently, the only element of $\Gamma$ mapping every element of $\mcA$ to itself is the identity).
\item For every $\gamma \in \Gamma$, $\{x : \gamma(x)= x \} \in \mcA$.
\item For every $g \in \Homeo(X)$, if there exists a partition $X=\bigsqcup_{i=1}^n A_i$ with $A_i \in \mcA$ such that $g$ coincides on each $A_i$ with some $\gamma_i \in \Gamma$, then $g \in \Gamma$.
\end{itemize}
We sometimes denote by $\Gamma_\mcA$ the subgroup $e_\mcA(\Gamma)$ of $\Aut(\mcA)$.

One says that the unit system $(\mcA,\Gamma)$ is \emph{finite} if $\mcA$ is finite, in which case $\Gamma$ is also finite. We say that a unit system $(\mcB,\Sigma)$ \emph{refines} another unit system $(\mcA,\Gamma)$ if $\mcB$ contains $\mcA$ and $\Sigma$ contains $\Gamma$.

Krieger~\cite[Lemma~2.1]{Krieger1979} proved that for any ample group there exists a sequence $(\mcA_n,\Gamma_n)_n$ of finite unit systems such that $(\mcA_{n+1},\Gamma_{n+1})$ refines $(\mcA_n,\Gamma_n)$ for all $n$ and
\[
\bigcup_n \mcA_n=\Clopen(X); \quad \bigcup_n \Gamma_n=\Gamma.
\]
We say that $(\mcA_n,\Gamma_n)$ is an \emph{exhaustive sequence} for $(X,\Gamma)$.

\emph{From now on (throughout the paper) the letter $X$ stands for the Cantor space}.

An action $\Gamma \actson X$ by homeomorphisms is \emph{minimal} if all of its orbits are dense; we say that $\varphi \in \Homeo(X)$ is minimal if the $\Z$-action $n \cdot x=\varphi^n(x)$ is minimal. We denote by $F(\varphi)$ the full group generated by $\{\varphi^n : n \in \Z\}$ (it is often denoted $\llbracket \varphi \rrbracket$ in the literature). This is a countable group which acts minimally on $X$; for any $x_0 \in X$, the subgroup
\[
F_{x_0}(\varphi)= \left\{ \gamma \in F(\varphi) : \gamma(O^+(x_0))=O^+(x_0)\right\}
\]
is an ample group (here $O^+(x_0)=\{ \varphi^n(x_0) : n \ge 0\}$). The orbits of $F_{x_0}(\varphi) \actson X$ are the same as the $\varphi$-orbits, except for the orbit $O(x_0)$ which splits into its positive part $O^+(x_0)$ and negative part $O^-(x_0)$. It is proved in~\cite{Melleray2023} that all minimal ample groups over $X$ can be realized in this way, so minimal ample groups and topological full groups of minimal $\Z$-actions are closely related.

Any action $\Gamma \actson X$ also induces an equivalence relation $\sim_\Gamma$ on $\Clopen(X)$, where $A \sim_\Gamma B$ iff there exists $\gamma \in \Gamma$ such that $\gamma(A)=B$. If $\Gamma$ is ample, then $\sim_\Gamma$ is \emph{full} in the following sense: whenever $A, B \in \Clopen(X)$ are such that $A= \bigsqcup_{i=1}^n A_i$, $B=\bigsqcup_{i=1}^n B_i$ and $A_i \sim_\Gamma B_i$ for all $i$ then $A \sim_\Gamma B$.


\begin{defi}
We say that two ample subgroups $\Gamma$, $\Lambda$ \emph{induce isomorphic relations on $\Clopen(X)$} if there exists $h \in \Homeo(X)$ such that for any clopen $A,B$ one has $(A \sim_\Gamma B) \Leftrightarrow (h(A) \sim_\Lambda h(B))$.
\end{defi}

Krieger's theorem alluded to in the introduction (and which we strengthen in Section~\ref{s:Krieger}) implies that if two ample groups $\Gamma$, $\Lambda$ induce isomorphic relations on $\Clopen(X)$ then $\Gamma$ and $\Lambda$ are conjugated in $\Homeo(X)$. In particular $R_\Gamma$ and $R_\Lambda$ are then orbit equivalent.


The following lemma, whose analogue for $\Z$-actions is due to Glasner and Weiss \cite{Glasner1995a}, is crucial for our approach.


\begin{lemm}[Glasner--Weiss]\label{l:Glasner_Weiss}
Assume that $\Gamma \le \Homeo(X)$ is an ample group. Fix $A,B \in \Clopen(X)$. If $\mu(A) < \mu(B)$ for all $\mu \in M(\Gamma)$ then there exists $\gamma \in \Gamma$ such that $\gamma(A)\subset B$.
\end{lemm}



We will use another property of $\sim_\Gamma$ (which is implicitly used in~\cite{Melleray2023}).

\begin{lemm}\label{l:equidecomposability_on_clopens_is_induced_by_full_group}
Let $\Gamma$ be an ample group over $X$, and $A$, $B$ be clopen subsets of~$X$.

Assume that $A \sim_\Gamma B$ and that $A_1 \subseteq A$, $B_1 \subseteq B$ are clopen sets such that $A_1 \sim_\Gamma B_1$. Then $A \setminus A_1 \sim_\Gamma B \setminus B_1$.
\end{lemm}

\begin{proof}
Fix an exhaustive sequence $(\mcA_n,\Gamma_n)$ of finite unit systems for $\Gamma$. By assumption, there exists $n$ and $\gamma$, $\delta \in \Gamma_n$ such that $A$, $A_1$, $B$, $B_1$ all belong to $\mcA_n$, $\gamma A=B$ and $\delta A_1=B_1$.

Given an atom $U$ of $\mcA_n$ and $C \in \mcA_n$ define
\[
n_C(U)= \left|\{V \in \Gamma_n U : V \subseteq C\}\right|.
\]

Since $A,B \in \mcA_n$, the existence of $\gamma \in \Gamma_n$ such that $\gamma A=B$ amounts to saying that $n_A(U)=n_B(U)$ for every atom $U$ of $\mcA_n$; similarly, we obtain $n_{A_1}(U)= n_{B_1}(U)$ for every $U$. Then for every $U$ we have
\[
n_{A \setminus A_1}(U)=n_A(U) - n_{A_1}(U) = n_B(U)-n_{B_1}(U) = n_{B \setminus B_1}(U)
\]
and that proves Lemma~\ref{l:equidecomposability_on_clopens_is_induced_by_full_group}.
\end{proof}


\section{Malleable subsets}

Throughout this section we fix an ample $\Gamma$ group over $X$ which acts minimally. Below, we will use without mention the fact that for any closed subset $K$ of $X$, any clopen subset of $K$ is the intersection of $K$ with a clopen subset of $X$.

\begin{defi}
[{see~\cite[Definition~4.11]{Giordano2004}}]
Let $K$ be a closed subset of~$X$. We say that $K$ is $\Gamma$-\emph{thin} if $\mu(K)=0$ for every $\mu \in M(\Gamma)$.
\end{defi}

The next lemma is standard.

\begin{lemm}\label{l:small_clopen}
Assume that $K$ is a closed, $\Gamma$-thin subset of $X$. Then for every $\varepsilon >0$ there exists a clopen subset $U$ such that $K \subset U$ and $\mu(U) \le \varepsilon$ for all $\mu \in M(\Gamma)$.
\end{lemm}

\begin{proof}
Pick a decreasing sequence $(U_n)_{n \in \N}$ of clopen subsets such that\linebreak ${\bigcap_n U_n = K}$. Given $n \in \N$, the function $\Phi_n \colon \mu \mapsto \mu(U_n)$ is continuous on $M(\Gamma)$ (endowed with its usual compact topology).

Furthermore, for every $\mu \in M(\Gamma)$ the sequence $(\Phi_n(\mu))_n$ decreases to $\mu(K)=0$. Since $M(\Gamma)$ is compact we can apply Dini's theorem and conclude that there is $n$ such that $\Phi_n(\mu) \le \varepsilon$ for all $\mu \in M(\Gamma)$, equivalently $\mu(U_n) \le \varepsilon$ for all $\mu \in M(\Gamma)$.
\end{proof}

Then the argument used to prove~\cite[Lemma~3.13]{Melleray2023} gives the following result.

\begin{lemm}\label{l:compatible_clopen}
Assume that $K$ is a closed $\Gamma$-thin subset. Let $U$ be a nontrivial clopen subset of $X$, and $A$ a clopen subset of $K$. Let $V \in \Clopen(X)$ be such that $A \subset V \cap K$ and $\mu(U)< \mu(V)$ for all $\mu \in M(\Gamma)$.

Then there exists $U' \in \Clopen(X)$ such that $U' \cap K =A$, $U' \subset V$, and $U' \sim_\Gamma U$.
\end{lemm}

\begin{proof}
By compactness of $M(\Gamma)$ (and minimality of the action) there exists $\eta >0$ such that both $\eta < \mu(U)$ and $\mu(U)+ \eta < \mu(V)$ for all $\mu \in M(\Gamma)$.

By Lemma~\ref{l:small_clopen}, there exists a clopen subset $W$ containing $K$ such that $\mu(W)< \eta$ for all $\mu \in M(\Gamma)$. Since $A$ is clopen in $K$, we have $A=K \cap B$ with $B$ clopen in $X$ and contained in $W \cap V$. Since $\mu(B) < \eta < \mu(U)$ for all $\mu \in M(\Gamma)$, we may apply Lemma~\ref{l:Glasner_Weiss} and conclude that there exists a clopen $C \subseteq U$ such that $C \sim_\Gamma B$.

For all $\mu \in M(\Gamma)$ we have
\[
\mu(U \setminus C) +\eta = \mu(U) - \mu(C) + \eta < \mu(V)-\mu(B) = \mu(V \setminus B)
\]
Applying Lemma~\ref{l:small_clopen} again, we pick a clopen subset $D$ of $V \setminus B$ containing $(K \cap V) \setminus A$ and such that $\mu(D)< \eta$ for all $\mu \in M(\Gamma)$. We then have $\mu(U \setminus C) < \mu(V \setminus (B \cup D))$ so by Lemma~\ref{l:Glasner_Weiss} there exists a clopen $E \subset V \setminus (B \cup D)$ and such that $E \sim_\Gamma U\setminus C$.

We conclude by setting $U'=B \cup E$.
\end{proof}

The way we used Lemma~\ref{l:Glasner_Weiss} above is typical of our arguments; from now on\linebreak we will not mention this lemma explicitly but it is very often in the background.

\begin{defi}
Let $K$ be a closed subset of $X$. We say that $K$ is $\Gamma$-\emph{étale} if for any $\gamma \in \Gamma$ and any clopen subset $A$ of $K$ the set $\gamma A \cap K$ is clopen in $K$.

We say that $K$ is $\Gamma$-\emph{malleable} if it is a closed, $\Gamma$-thin and $\Gamma$-étale subset of $X$.
\end{defi}

The terminology ``étale'' comes from work of Giordano--Putnam--Skau which itself originates from operator algebra theory. Our definition amounts to saying that the restriction of $R_\Gamma$ to $K$ is étale for the topology induced by the topology of~$R_\Gamma$, that is, $K$ is $R_\Gamma$-étale; see~\cite[Definition~2.1]{Giordano2004}. In~\cite{Giordano2004}, one requires compatibility between $K$ and the topology of the étale equivalence relation under consideration; here, analogously, we ask for compatibility with the acting ample group. Our malleable sets correspond to the $R$-closed, $R$-étale and $R$-thin subsets considered in~\cite{Giordano2004} and later in~\cite{Giordano2008} and~\cite{Matui2008}.

Assume that $K$ is $\Gamma$-étale. Given an involution $\gamma \in \Gamma$, define $\gamma_K \colon K \to K$ by setting $\gamma_K(x)= \gamma(x)$ if $\gamma(x) \in K$, $\gamma_K(x)=x$ otherwise. Denote $K_1= \{x \in K : \gamma(x) \in K\}$. Then $K_1$ is clopen in $K$ since $K$ is $\Gamma$-étale, and $\gamma_K(K_1)= K_1$. From this we obtain that $\gamma_K$ is a homeomorphic involution of $K$.

\begin{defi}
Let $K$ be a $\Gamma$-étale subset of $X$, and $\gamma \in \Gamma$ an involution. We say that $\gamma$ is $K$-\emph{compatible} if $\gamma(K)=K$.
\end{defi}

\begin{lemm}
Let $K$ be a $\Gamma$-étale subset of $X$, and $\gamma \in \Gamma$ be an involution. Then there exists a $K$-compatible involution $\delta \in \Gamma$ such that $\delta_K=\gamma_K$.
\end{lemm}

\begin{proof}
The sets $\{x \in K : \gamma(x) \in K \}$ and $\{x \in K : \gamma(x) \ne x\}$ are both clopen in $K$; denote by $L$ their intersection. Since $\gamma$ is an involution we have $\gamma(L)=L$.

Choose a clopen subset $U$ of $X$ which contains $L$ and is disjoint from $K\setminus L$, and set $V=U \cap \gamma(U)$. Then $V$ is clopen, $V \cap K=L$ and $V= \gamma (V)$. We can then define $\delta \in \Gamma$ by setting $\delta(x)=\gamma(x)$ for all $x \in V$ and $\delta(x)=x$ for all $x \not \in V$.

By construction, $\delta(x)=\gamma(x)$ for all $x \in L$, and $\delta(x)=x$ for all $x \in K \setminus L$, so that $\delta_K=\gamma_K$.
\end{proof}

\begin{defi}
Let $\Gamma$ be a minimal ample subgroup of $X$ and $K$ a $\Gamma$-étale subset. Let $(\mcB,\Delta)$ be a finite unit system with $\Delta \le \Gamma$.

We say that $(\mcB,\Delta)$ is \emph{$K$-compatible} if for all atoms $A$, $B$ of $\mcB$ and every $\delta \in \Delta$ such that $\delta(A)=B$, if $K \cap A$ and $K \cap B$ are both nonempty then we have $\delta(K \cap A)= K \cap B$.
\end{defi}

This extends a definition given in~\cite[Definition~3.10]{Melleray2023}, where we only considered closed sets $K$ such that for all $x\ne y \in K$ one has $y \not \in \Gamma x$, which we called $\Gamma$-sparse sets. The assumption of $K$-compatibility amounts to saying that, if $A,B$ are two elements of the same $\Delta$-orbit that both intersect $K$, then the involution in $\Delta$ which maps $A$ to $B$ and is equal to the identity outside of $A \cup B$ is $K$-compatible.

\begin{lemm}\label{l:refining_to_compatible}
Assume that $K$ is $\Gamma$-étale and let $(\mcA,\Delta)$ be a finite unit system with $\Delta \le \Gamma$. Then there exists $\mcA'$ refining $\mcA$ and $\Delta'$ such that $(\mcA',\Delta')$ is\linebreak a $K$-compatible finite unit system and $\Delta \le \Delta' \le \Gamma$.
\end{lemm}

Since any two finite unit systems $(\mcA,\Delta)$ and $(\mcB,\Lambda)$ with $\Delta, \Lambda \le \Gamma$ have a common refinement $(\mcC,\Sigma)$ with $\Sigma \le \Gamma$, it follows from Lemma~\ref{l:refining_to_compatible} and the existence of an exhaustive sequence of finite unit systems for $\Gamma$ that if $K$ is $\Gamma$-étale then there exists an exhaustive sequence of $K$-compatible finite unit systems for $\Gamma$.

\begin{proof}
The argument proceeds by cutting the $\Delta$-orbit of each atom of $\mcA$. First, fix some such orbit $\tau$.

Let $U \ne V$ be two elements of $\tau$ and $\delta \in \Delta$ the involution such that $\delta U=V$ and $\delta(x)=x$ for every $x \not \in U \sqcup V$. We first consider $K_1(\delta) = \{x \in K \cap U : \delta(x) \in K\}$, which is clopen in $K$; we pick a clopen set $U' \subseteq U$ such that $U' \cap K=K_1(\delta)$. Then $\delta(K_1(\delta))=\delta(U' \cap K)$ is clopen in $K$ and contained in $V$, so we can find a clopen subset $V'$ of $V$ such that $V' \cap K= \delta(K_1(\delta))$. Set $U_1 = U' \cap \delta V'$, $V_1= \delta (U_1)$.

By definition, $K_1(\delta)= U_1 \cap K$ and $\delta(U_1 \cap K)= \delta(K_1(\delta))=V_1 \cap K$.

If $(U \setminus U_1) \cap K= \emptyset$ or $(V \setminus V_1) \cap K= \emptyset$ we let $U_2=U \setminus U_1$, $V_2=\delta(U_2)$ and stop partitioning $U$, $V$. Else, let $U''=U \setminus U_1$, $V''=V \setminus V_1$. Then $\delta(U'' \cap K) \cap K=\emptyset$. Since these two sets are closed, we can find a clopen set $U_2 \subset U''$ containing $U'' \cap K$ and such that $\delta(U_2) \cap K= \emptyset$, and set $V_2= \delta(U_2)$. Finally we define $U_3= U \setminus (U_1 \sqcup U_2)$ and $V_3=\delta(U_3)$.

We now have $\delta(U_1 \cap K)= V_1 \cap K$; $V_2=\delta(U_2)$ and $V_2 \cap K=\emptyset$; $V_3=\delta(U_3)$ and $U_3 \cap K=\emptyset$.

We do this for every pair of distinct $U,V$ in $\tau$ and choose a finite $\Delta$-invariant partition of $\tau$ refining all those partitions. We obtained our desired refinement of $\tau$. Doing this for every $\tau$ gives us $\mcA'$. We then let $\Delta'$ be the group of all homeomorphisms of $X$ which preserve $\mcA'$ and coincide on each atom $U$ of $\mcA$ with some $\delta_U \in \Delta$.
\end{proof}


Given a $\Gamma$-étale $K \subset X$, we denote by $\Gamma_K$ the smallest full subgroup of $\Homeo(K)$ which contains $\gamma_K$ for every involution $\gamma \in \Gamma$. For every $x \in X$ and every $\gamma \in \Gamma$ there exists an involution $\delta \in \Gamma$ such that $\delta(x)=\gamma(x)$; hence the restriction of $R_\Gamma$ to some malleable $K$ is induced by $\Gamma_K$.

The group $\Gamma_K$ can equivalently be described as follows.

\begin{lemm}\label{l:induced_group_is_what_you_expect}
%
\!\!
Assume that $K$ is $\Gamma$-étale. Then $\Gamma_K\Mk =\Mk\{\gamma_{|K}\Mk :\Mk \gamma \in \Gamma \text{ and } \gamma(K)\Mk=\Mk K\}$.\linebreak In particular, $\Gamma_K$ is an ample group over $K$.
\end{lemm}

\begin{proof}
Fix an exhaustive sequence $(\mcA_n,\Gamma_n)$ of $K$-compatible finite unit systems for $\Gamma$ (see the remark following the statement of Lemma~\ref{l:refining_to_compatible}).

Pick $g \in \Gamma_K$. There exists a clopen partition $(U_i)_{i=1,\dots,p}$ of $K$ such that for each $i$ the restriction of $g$ to $U_i$ coincides on $U_i$ with a product of $K$-compatible involutions. Let $A$ be a finite set of $K$-compatible involutions witnessing that fact, and define $\Delta=\langle A\rangle$ (a finite subgroup of $\Gamma$); note that $\delta(K)=K$ for all $\delta \in \Delta$, in particular restrictions of elements of $\Delta$ to $K$ form a finite subgroup of $\Homeo(K)$.

Next, pick a finite subalgebra $\mcB$ of $\Clopen(K)$ containing every $U_i$ and such that $\delta(B) \in \mcB$ for all $\delta \in \Delta$. Choose a Boolean subalgebra $\tilde \mcB$ of $\Clopen(X)$ such that $\{ \tilde B \cap K : \tilde B \in \tilde \mcB \}=\mcB$ and $\tilde \mcB$ is $\Delta$-invariant. For $n$ large enough we have that $\mcA_n$ refines $\tilde \mcB$ and $\Gamma_n$ contains $\Delta$. For every atom $U$ of $\mcA_n$ which intersects $K$, there exists $\delta \in \Delta$ such that $\delta$ and $g$ coincide on $U \cap K$. We must then have $\delta(U \cap K)= g(U \cap K)$, and any two such $\delta$ coincide on $U$ since $\Delta$ is contained in $\Gamma_n$ and $(\mcA_n,\Gamma_n)$ is a unit system.

Denoting this $\delta$ by $\delta_U$, we define $\gamma \in \Homeo(X)$ by setting $\gamma(x)=\delta_U(x)$ for every atom $U$ of $\mcA_n$ which intersects $K$ and every $x \in U$, $\gamma(x)=x$ elsewhere. By definition we have both that $\gamma \in \Gamma$ and $\gamma_{|K}=g$.

Conversely, assume that $\gamma \in \Gamma$ is such that $\gamma(K)=K$. Then find $n \in \N$ such that $\gamma \in \Gamma_n$. Since $(\mcA_n,\Gamma_n)$ is $K$-compatible, for each atom $U$ of $\mcA_n$ such that $U \cap K \ne \emptyset$ there exists a unique involution $\gamma_U$ in $\Gamma_n$ such that $\gamma_U(U \cap K)= \gamma(U \cap K)$ and $\gamma_U(x)=x$ for all $x \not \in U \cup \gamma(U)$. By definition of a unit system, $\gamma$ and $\gamma_U$ coincide on $U$, and $\gamma_U$ is $K$-compatible. Let $\lambda$ be the homeomorphism of $K$ which coincides with $\gamma_U$ on each $U \cap K$ whenever $U \cap K \ne \emptyset$ and coincides with the identity everywhere else. Then $\lambda \in \Gamma_K$ and $\lambda=\gamma_{|K}$.
\end{proof}

\begin{defi}
Assume that $K$ is $\Gamma$-étale and that $(\mcA,\Lambda)$ is a $K$-compatible finite unit system with $\Lambda \le \Gamma$. We let $\mcA_K$ denote the Boolean subalgebra of $\Clopen(K)$ induced by $\mcA$, and $\Lambda_K=\{\lambda_{|K} : \lambda\in \Lambda \text{ and } \lambda(K)=K\}$.

Note that by $K$-compatibility we have that $(\mcA_K,\Lambda_K)$ is a finite unit system in $K$.
\end{defi}

\begin{lemm}
Assume that $(\mcA_n,\Gamma_n)$ is an exhaustive sequence of $K$-compatible finite unit systems for $\Gamma$. Then $(\mcA_{n,K},\Gamma_{n,K})$ is an exhaustive sequence of finite unit systems for $\Gamma_K$.
\end{lemm}

\begin{proof}
Clearly $\bigcup_n \mcA_{n,K}=\Clopen(K)$. Given $\sigma \in \Gamma_K$, by Lemma~\ref{l:induced_group_is_what_you_expect} there exists $\gamma \in \Gamma$ such that $\sigma=\gamma_{|K}$, whence $\sigma \in \Gamma_{n,K}$ for $n$ large enough.
\end{proof}

\setcounter{section}{4}

\section{A first absorption theorem}
\begin{defi}
Let $K$ be a closed subset of $X$, and $\Delta$ be a subgroup of $\Homeo(K)$. We denote by $R_\Gamma(K,\Delta)$ the finest equivalence relation $S$ which is coarser than $R_\Gamma$ and is such that $(x,\delta x) \in S$ for all $x \in K$ and all $\delta \in \Delta$.
\end{defi}

Informally, $R_\Gamma(K,\Delta)$ is obtained from $R_\Gamma$ by joining the $\Gamma$-orbits of $\Delta$-equivalent elements of $K$, and leaving untouched the $\Gamma$-orbits which do not intersect $K$.

\setcounter{section}{5}

\section{The absorption theorem}

The argument in the previous section could be pushed further (for instance, one can recover a proof of~\cite[Theorem~6.4]{Melleray2023}) but we want to establish a stronger absorption theorem, and that seems to require some significant extra work.

We again fix an ample group $\Gamma$ over $X$ which acts minimally. We need to recall some more vocabulary from~\cite{Melleray2023}.

\begin{defi}
We say that a clopen partition $\mcA=(A_{i,j})_{(i,j) \in I}$ is a \emph{$\sim_\Gamma$-partition} if:
\[
\forall i,j,k \quad ((i,j) \in I \quad\text{and}\quad (i,k) \in I) \To (A_{i,j} \sim_\Gamma A_{i,k})
\]
Denoting by $I_i=\{j : (i,j) \in I\}$ we say that $\{A_{i,j} : j \in I_i\}$ is an \emph{$\mcA$-orbit}.

A \emph{fragment} of an $\mcA$-orbit $(A_{i,j})_{j \in I_i}$ is a family $(B_{i,j})_{j \in I_i}$ of nonempty clopen subsets of $X$ such that $B_{i,j} \subseteq A_{i,j}$ for all $j \in I_i$ and $B_{i,k} \sim_\Gamma B_{i,j}$ for all $j,k \in I_i$.
\end{defi}

For the reader who is already familiar with Kakutani--Rokhlin partitions and cutting-and-stacking arguments, orbits (and their fragments) are unordered counterparts of Kakutani--Rokhlin partitions.
A fragment of an orbit is what one obtains by applying a cutting procedure to build a finer partition than what one started with. Note that whenever $(\mcA,\Delta)$ is a finite unit system with $\Delta \le \Gamma$, one can view $\mcA$ as\linebreak a $\sim_\Gamma$-partition by grouping together atoms of $\mcA$ which belong to the same $\Delta$-orbit.

We say that a $\sim_\Gamma$-partition $\mcB$ \emph{refines} another $\sim_\Gamma$-partition $\mcA$ if every $\mcB$-orbit can be written as a disjoint union of fragments of $\mcA$-orbits (intuitively, $\mcB$ has been obtained from $\mcA$ by cutting some $\mcA$-orbits, then grouping some fragments together). We note that we sometimes identify a clopen partition and the Boolean algebra it generates, which should cause no risk of confusion.



The following lemma is the key step of our proof of the absorption theorem (this is what allows the ``Hilbert--Bratteli hotel'' argument to go through). The idea is the following: in order to prove that $R_\Gamma$ can absorb some small extension of itself, we want to show that $R_\Gamma$ is itself obtained by repeating $\omega$ times the same small extension; first and foremost, one needs to show that $R_\Gamma$ can by realized as a prescribed small extension of $R_\Lambda$ for some minimal ample group $\Lambda$.

\begin{lemm}\label{l:everyone_is_a_small_extension}
Let $Y$ be a compact, $0$-dimensional metric space and $\Delta \le \Sigma$ two ample groups over $Y$.

There exists a minimal ample group $\Lambda \le \Gamma$, a closed subset $K$ which is both $\Lambda$- and $\Gamma$-malleable, and a homeomorphism $h \colon Y \to K$ such that $h \Delta h^{-1} = \Lambda_K$, $h \Sigma h^{-1} = \Gamma_K $ and $R_\Gamma=R_\Lambda(K,\Gamma_K)$.
\end{lemm}

\begin{proof}
We fix an exhaustive sequence $(\mcA_n,\Sigma_n)$ of finite unit systems for $(Y,\Sigma)$ and let $\Delta_n=\Sigma_n \cap \Delta$.

Denote by $h_n$ the number of atoms of $\mcA_n$. We pick an exhaustive sequence $(\mcB_n,\Gamma_n)$ of finite unit systems for $(X,\Gamma)$; for all $n$ we view $\mcB_n$ as a $\sim_\Gamma$-partition so that the $\mcB_n$- and $\Gamma_n$-orbits of every atom of $\mcB_n$ coincide.

We also assume that for any $n$ there are more than $h_n$ distinct $\mcB_n$-orbits (which may be achieved by cutting them if necessary). Denoting by $k_n$ the number of atoms of $\mcB_n$, we may also ensure (because $\Gamma$ acts minimally) that every $\mcB_{n+1}$-orbit contains more than $(k_n+1)h_{n+1}$ disjoint fragments of every $\mcB_n$-orbit and that $k_{n} \ge n h_{n+1}$.

To initialize the construction, we add the assumption that $\mcA_0$, $\mcB_0$, $\Sigma_0$ and $\Gamma_0$ are all trivial.


\begin{step}[{Building a copy of $Y$}]\label{step1}
The first step consists of choosing, within each~$\mcB_n$, some atoms which form a finer and finer approximation $K$ of $Y$; and doing so in such a way that these atoms generate a sequence of clopen partitions of $K$ which is induced by a copy of $\Sigma$. We also want to introduce ``cuts'' into the $\mcB_n$-orbits so as to mimic the way $\Sigma$-orbits of $\mcA_n$ are cut when one moves down from $\Sigma$ to $\Delta$.


By first choosing some elements of $\mcB_n$, then grouping them together with other elements of $\mcB_n$, we build a sequence of maps $\Phi_n \colon \mcA_n \to \mcB_n$ and a sequence of equivalence relations $\sim_n$ on $\mcB_n$; intuitively, $\Phi_n(\mcA_n)$ is our approximation (at step $n$) of the domain of the copy of $(Y,\Delta,\Sigma)$ that we are trying to build, and the restriction of $\sim_n$ to $\Phi_n(\mcA_n)$ mimics $\sim_{\Delta_n}$.

\pagebreak
Explicitly, our construction proceeds by enforcing the following conditions:
\begin{enumerate}
\item $\Phi_n \colon \mcA_n \to \mcB_n$ is injective, and for any $A,B \in \mcA_n$ we have
\[
\left(\Phi_n(A) \sim_{\Gamma_n} \Phi_n(B)\right) \Leftrightarrow (A \sim_{\Sigma_n} B) \, ; \ \left(\Phi_n(A) \sim_n \Phi_n(B)\right) \Leftrightarrow (A \sim_{\Delta_n} B)
\]
\item For any $U,V \in \mcB_n$, $U \sim_n V \To U \sim_{\Gamma_n} V$.

\item If $U \in \mcB_{n}$ is such that the $\mcB_n$-orbit $\tau_U$ of $U$ does not intersect the image of~$\Phi_n$, then all elements of $\tau_U$ are $\sim_{n}$-equivalent. Else, every element of $\tau_U$ is $\sim_{n}$-equivalent to some $\Phi_{n}(A)$ for $A \in \mcA_{n}$.
\item If $A \in \mcA_n$, $B \in \mcA_{n+1}$ are such that $B \subseteq A$ then $\Phi_{n+1}(B) \subseteq \Phi_n(A)$.

Denote by $\mcC_n$ the $\sim_\Gamma$-partition whose atoms are the same as those of $\mcB_n$ and whose orbits are the $\sim_n$-classes. Each $\mcB_n$-orbit is a union of $\mcC_n$-orbits.

\item Given $A \in \mcA_{n+1}$, the $\sim_{n+1}$-class of $\Phi_{n+1}(A)$ is obtained by joining one distinguished fragment of the $\mcC_{n}$-orbit of every $\Phi_{n+1}(A')$ for $A' \in \Delta_{n+1}A$ and fragments of the other $\mcB_n$-orbits, so that a fragment of every $\mcB_n$-orbit appears at least once. In particular, $\mcC_{n+1}$ refines $\mcC_n$ and each $\mcC_{n+1}$-orbit contains a fragment of every $\mcC_n$-orbit.
\end{enumerate}

To see that this is indeed possible, assume that this construction has been carried out up to some $n$ (for $n=0$ there is nothing to do). Below, for $A$ an atom of $\mcA_{n+1}$ we denote by $A'$ the unique atom of $\mcA_n$ which contains $A$.

We first define an injective $\Phi_{n+1} \colon \mcA_{n+1} \to \mcB_{n+1}$. For any $A \in \mcA_{n+1}$ we ask that $\Phi_{n+1}(A) \subseteq \Phi_n(A')$, and require also that, for any $A,B \in \mcA_{n+1}$, $\Phi_{n+1}(A)$ and $\Phi_{n+1}(B)$ belong to the same $\mcB_{n+1}$-orbit iff $A$ and $B$ belong to the same $\Sigma_{n+1}$-orbit. This is feasible because the number of $\mcB_{n+1}$-orbits is greater than $h_{n+1}$ (hence larger than the number of $\Sigma_{n+1}$-orbits of atoms of $\mcA_{n+1}$) and, for any $A' \in \mcA_n$, in any $\mcB_{n+1}$-orbit there are more than $h_{n+1}$ atoms contained in $A'$, so that we can guarantee the injectivity of $\Phi_{n+1}$.

Next, fix a $\mcB_{n+1}$-orbit $\tau$ which intersects $\Phi_{n+1}(\mcA_{n+1})$. We start building $\sim_{n+1}$ on~$\tau$, by choosing for each $A \in \mcA_{n+1}$ such that $\Phi_{n+1}(A) \in \tau$ a fragment of the $\mcC_n$-orbit of $\Phi_n(A')$ which contains $\Phi_{n+1}(A)$ and is contained in $\tau$. We do this in such a way that $\Phi_{n+1}(A)$ and $\Phi_{n+1}(B)$ belong to the same fragment iff $A' \in \Delta_{n} B'$. Again, this is feasible because $\tau$ contains many fragments of every $\mcB_n$-orbit, hence of every $\mcC_n$-orbit. Then we include in the $\sim_{n+1}$-class of $\Phi_{n+1}(A)$ the union of the fragments we have chosen for all $B \in \Delta_{n+1}A$ (and nothing else at this point).

The union of fragments of $\mcC_n$-orbits on which we have so far defined $\sim_{n+1}$ in~$\tau$ is actually a union of fragments of $\mcB_n$-orbits: indeed, for each $\mcB_n$-orbit $\rho$ which intersects $\Phi_n(\mcA_n)$ we have selected whole fragments, because $\rho$ is the union of $\sim_n$-classes of all $\Phi_n(A')$ contained in $\rho$ (and the fragments of other $\mcB_n$-orbits have not yet been taken under consideration). Thus, at this point, each partial $\sim_{n+1}$-equivalence class in $\tau$ consists of a union of a most $h_{n+1}$ fragments of $\mcB_{n}$-orbits. Inside $\tau$, we have at least $(k_n+1)h_{n+1}$ such fragments of each $\mcB_n$-orbit at our disposal, hence we can distribute the remaining fragments among the (at most $h_{n+1}$) disjoint partial $\sim_{n+1}$-classes contained in $\tau$ in such a way that each $\sim_{n+1}$-class contains at least one fragment of every $\mcB_n$-orbit, and every atom of $\tau$ is in the $\sim_{n+1}$-class of some $\Phi_{n+1}(A)$.

Performing the above procedure for every $\tau$ which intersects $\Phi_{n+1}(\mcA_{n+1})$, we obtain the desired relation $\sim_{n+1}$.

For all $n \in \N$ we let $K_n= \bigsqcup_{A \in \mcA_n} \Phi_n(A)$. This is a clopen set, $K_{n+1} \subseteq K_n$ for all $n$ and we define $K=\bigcap_{n \in \N} K_n$.

The sequence $(\Phi_n)_n$ induces an isomorphism $\Phi \colon \Clopen(Y) \to \Clopen(K)$, and we let $h \colon Y \to K$ be the corresponding homeomorphism.
\end{step}
%\medskip
\begin{step}[{Defining $\Lambda$}]\label{step2}
We now want to use our previous work in order to define a copy $\Lambda$ of $\Delta$.

Denote by $\Lambda_n$ the group of permutations of the atoms of $\mcB_n$ which map each $\sim_n$-class to itself. We define embeddings $i_n \colon \Lambda_n \to \Lambda_{n+1}$ so that (among other properties) for any $\lambda \in \Lambda_n$ and any $B \in \mcB_n$, if $B_1,\dots,B_p$ are the atoms of $\mcB_{n+1}$ such that $B= \bigsqcup_{i=1}^p B_i$ then $\lambda (B)= \bigsqcup_{i=1}^p i_n(\lambda)(B_i)$.

To explain how $i_n$ is defined, we fix a $\sim_n$-orbit $\tau$ and assume first that there is $A \in \mcA_n$ such that $\Phi_n(A) \in \tau$. Given another atom $U$ of $\tau$, denote $\lambda_{U}$ the unique involution in $\Lambda_n$ which maps $\Phi_n(A)$ to $U$ and leaves all other elements of $\mcB_n$ fixed.

First, we consider the case where $U=\Phi_n(\delta A)$ with $\delta \in \Delta_n$. We can write $A=\bigsqcup_{i=1}^p A_i$ with $A_i \in \mcA_{n+1}$ and set
\[
i_n(\lambda_U)(\Phi_{n+1}(A_i))=\Phi_{n+1}(\delta A_i)
\]
We still have to define $i_n(\lambda_U)$ on some atoms of $\mcB_{n+1}$ contained in $\Phi_n(A),\Phi_n(\delta A)$. In each $\mcC_{n+1}$-orbit there remain as many atoms contained in $\Phi_n(\delta A)$ where $i_n(\lambda_U)$ has not been defined as there are such atoms contained in $\Phi_n(A)$ (because $\mcC_{n+1}$ refines $\mcC_n$, every $\mcC_{n+1}$-orbit has as many atoms contained in $\Phi_n(A)$ and $\Phi_n(\delta A)$). We may thus match those atoms arbitrarily to define $i_n(\lambda_U)$ so that it maps each of these atoms to an atom in the same $\sim_{n+1}$-class.

The second case is when $U$ is not in the image of $\Phi_n(\mcA_n)$. Again, each $\mcC_{n+1}$-orbit has as many atoms contained in $U$ and in $\Phi_n(A)$, and we match those arbitrarily to define $i_n(\lambda_U)$.

If our $\sim_n$-orbit $\tau$ does not intersect the image of $\Phi_n$, then we choose $V \in \tau$ at random, define $\lambda_U$ as above (for $U$ another element of $\tau$) and define $i_n(\lambda_U)$ as in the previous paragraph (any involution which extends $\lambda_U$ and fixes setwise each $\sim_{n+1}$-class will do the job).

Once all this is done, we have completely defined $i_n \colon \Lambda_n \to \Lambda_{n+1}$.

The inductive limit $\Lambda$ of the sequence $(\Lambda_n,i_n)_n$ naturally acts on $\Clopen(X)$, and we view it as a subgroup of $\Homeo(X)$. By construction, this group is ample.

Denoting also by $i_n$ the embedding of $\Lambda_n$ in $\Lambda$, the $i_n(\Lambda_n)$-orbit of every element of $\mcC_{n}$ coincides with its $\mcC_{n}$-orbit. Given a nonempty $U \in \Clopen(X)$, let $n$ be such that $U $ is a union of atoms of $\mcC_n$. Since every $\mcC_{n+1}$-orbit contains a fragment of every $\mcC_n$-orbit, hence an atom contained in $U$, we have $i_{n+1}(\Lambda_{n+1})U=X$. We conclude that $\Lambda$ acts minimally on $X$.

In each $\mcC_{n+1}$-orbit there are at least $k_n$ elements, and at most $h_{n+1}$ of those belong to $K_{n+1}$; since $h_{n+1}/k_n \xrightarrow[n \to + \infty]{} 0$, we conclude that $\mu(K)=0$ for every $\mu \in M(\Lambda)$, so $K$ is $\Lambda$-thin.

By construction, if $A,B \in \mcA_n$ and $\lambda \in \Lambda_n$ are such that $\lambda \Phi_n(A)=\Phi_n(B)$ we can write $A= \bigsqcup_{i=1}^p A_i$, $B=\bigsqcup_{i=1}^p B_i$ with $A_i, B_i \in \mcA_{n+1}$ and $i_n(\lambda)(\Phi_{n+1}(A_i))=\Phi_{n+1}(B_i)$. By induction, it follows that $i_n(\lambda)(\Phi_n(A) \cap K)= \Phi_n(B) \cap K$. This implies that $K$ is $\Lambda$-étale. Finally, $K$ is $\Lambda$-malleable.

Our definition of $i_n$ also ensures that $h \Delta h^{-1}=\Lambda_{K}$ (where $h \colon Y \to K$ is the homeomorphism defined at the end of the first step).

Unfortunately we are not quite done yet: at this stage there is no reason why $\Lambda$ should be contained in $\Gamma$. To remedy this, we replace $\Gamma$ by a conjugate which we construct now.
\end{step}
%\medskip
\begin{step}[{Defining a conjugate of $\Gamma$}]
\label{step3}
We now want to build an ample group $\tilde \Gamma$ which contains $\Lambda$ and is such that one can write $\tilde \Gamma=\bigcup_n \tilde \Gamma_n$, so that the $\tilde \Gamma_n$-orbit of every element of $\mcB_n$ coincides with its $\mcB_n$-orbit. We also need to guarantee that $h \Sigma h^{-1} = \tilde{\Gamma}_K$. Again this group is obtained as an inductive limit of finite permutation groups.


For all $n$, we let $\tilde \Gamma_n$ denote the group of all permutations of $\mcB_n$ which map each $\mcB_n$-orbit to itself; note that $\Lambda_n$ is a subgroup of $\tilde \Gamma_n$. Similarly to what we did above, we build embeddings $j_n \colon \tilde \Gamma_n \to \tilde \Gamma_{n+1}$ to form an adequate inductive limit.

First, we ask that $j_n$ coincide with $i_n$ on $\Lambda_n$. Next, let $\tau$ be a $\mcB_n$-orbit; if $\tau$ is also a $\mcC_n$-orbit then we already know how to extend elements of $\tilde \Gamma_n$ whose support is contained in $\tau$ (those belong to $\Lambda_n$ so we simply apply $i_n$). So assume that $\tau$ intersects the image of $\Phi_n$ and let $A_1,\dots,A_p \in \mcA_n$ be such that $p \ge 2$ and $\tau$ is the disjoint union of the $\mcC_n$-orbits of $\Phi_n(A_1),\dots,\Phi_n(A_p)$.

The $\Sigma_n$-orbit of $A_1$ is the disjoint union of the $\Delta_n$-orbits of $A_1, \dots,A_p$. For $i \in \{2,\dots,p\}$ we let $\sigma_i$ be the involution in $\Sigma_n$ which maps $A_1$ to $A_i$ and fixes all other atoms, and denote $\gamma_i$ the element of $\tilde \Gamma_n$ which maps $\Phi_n(A_1)$ to $\Phi_n(A_i)$ and fixes all other atoms.

Fix $i \in \{2,\dots,p\}$. Write $A_1=\bigsqcup_{k=1}^q A_{1,k}$, $A_i=\bigsqcup_{k=1}^q \sigma_i(A_{1,k})$ with $A_{1,k} \in \mcA_{n+1}$, then set
\[
j_n(\gamma_i)(\Phi_{n+1}(A_{1,k}))=\Phi_{n+1}(\sigma_i A_{1,k})
\]

On each atom of $\mcB_{n+1}$ which does not intersect $\Phi_n(A_1) \sqcup \Phi_n(A_i)$ we must have $j_n(\gamma_i)$ coincide with the identity. Now, we observe that in each $\mcC_{n+1}$-orbit there are as many atoms contained in $\Phi_{n}(A_1)$ on which $j_n(\gamma_i)$ has not yet been defined as there are such atoms in $\Phi_n(A_i)$. This follows from the way $\sim_{n+1}$ has been defined above: so far we have defined what $j_n(\gamma_i)$ does to atoms contained in the $\mcC_n$-fragments intersecting $\Phi_{n+1}(\mcA_{n+1})$, and the remainder consists of unions of fragments of $\mcB_n$-orbits. We then match those atoms arbitrarily to define $j_n(\gamma_i)$ so that, on each atom $U$ which is not in $\Phi_{n+1}(\mcA_{n+1})$, $j_n(\gamma_i)U$ is $\sim_{n+1}$-equivalent to $U$ (this property is used at the end of the proof).

We apply this procedure to every $\mcB_n$-orbit on which $\sim_n$ is nontrivial. This finally defines $j_n$, and we obtain an ample group $\tilde \Gamma$ by considering the inductive limit of $(\tilde \Gamma_n,j_n)$. By construction we have $h \Sigma h^{-1}= \tilde{\Gamma}_K$. Denoting $j_n$ the embedding of $\tilde \Gamma_n$ into $\tilde \Gamma$, the $j_n(\tilde \Gamma_n)$-orbit of every element of $\mcB_n$ coincides with its $\Gamma_n$-orbit, and $(\mcB_n,j_n(\tilde \Gamma_n))_n$ is an exhaustive sequence of unit systems for $\tilde \Gamma$. It follows from Krieger's theorem that $\Gamma$ and $\tilde \Gamma$ are conjugate.
\end{step}

%\medskip
\begin{step}[{Checking that the construction satisfies the desired conditions}]\label{step4}
By definition, $\tilde \Gamma$ contains $\Lambda$. Since we already noted that $\mu(K)=0$ for all $\mu \in M(\Lambda)$, we also have $\mu(K)=0$ for all $\mu \in M(\tilde \Gamma)$.

Similarly to Step~\ref{step2}, the definition of each $j_n$ also ensures that for $A,B \in \mcA_n$ and $\gamma \in \tilde \Gamma_n$ such that $\gamma \Phi_n(A)=\Phi_n(B)$ we have $j_n(\gamma)(\Phi_n(A) \cap K)=\Phi_n(B) \cap K$. Hence $K$ is $\tilde \Gamma$-malleable. We have already noted that $h \Delta h^{-1} = \Lambda_K$ and $h \Sigma h^{-1} = \tilde \Gamma_K$.

Since $\Lambda \le \tilde \Gamma$ we have that $R_{\Lambda}(K,\tilde{\Gamma}_K)$ is contained in $R_{\tilde \Gamma}$. To see the converse inclusion, let $\gamma \in \tilde \Gamma$ and pick $n$ such that $\gamma \in j_n(\tilde \Gamma_n)$. Let $U$ be an atom of $\mcB_n$. If $\gamma U$ and $U$ belong to the same $\mcC_n$-orbit then $\gamma$ coincides on $U$ with an element of $\Lambda$, so $\gamma x \in \Lambda x$ for every $x \in U$. If $\gamma U$ and $U$ belong to different $\mcC_n$-orbits, then there exist $\lambda_1, \lambda_2 \in i_n(\Lambda_n)$, $A_1,A_2 \in \Phi_n(\mcA_n)$ and $\sigma \in \Sigma_n \setminus \Delta_n$ such that $\sigma(A_1)=A_2$ and $\lambda_1 U = \Phi_n(A_1)$, $\lambda_2 \Phi_n(A_2)=\gamma U$. Let $\gamma'$ be the involution of $\tilde \Gamma_n$ which maps $\Phi_n(A_1)$ to $\Phi_n(A_2)$ and fixes all other atoms. By construction, we have that $j_n(\gamma')(K)=K$, and for every $x \in X \setminus K$ we have $j_n(\gamma')(x) \in \Lambda x$. Since $\gamma$ coincides on $U$ with $i_n(\lambda_2)j_n(\gamma')i_n(\lambda_1)$, we see that $(x,\gamma x) \in R_\Lambda(K,\tilde{\Gamma}_K)$ for every $x \in U$. We finally conclude that $R_{\tilde \Gamma}=R_\Lambda(K,\tilde{\Gamma}_K)$.

We have obtained the desired result for $\tilde \Gamma$ instead of $\Gamma$; since $\Gamma$ and $\tilde \Gamma$ are conjugate in $\Homeo(X)$ this is enough.\qedhere
\end{step}\let\qed\relax
\end{proof}

\begin{thebibliography}{GMPS10}
\bibitem[GMPS08]{Giordano2008} Thierry Giordano, Hiroki Matui, Ian F. Putnam, and Christian F. Skau, \emph{The absorption theorem for affable equivalence relations}, Ergodic Theory Dyn. Syst. 28 (2008), no. 5, 1509--1531.
\bibitem[GMPS10]{Giordano2010} Thierry Giordano, Hiroki Matui, Ian F. Putnam, and Christian F. Skau, \emph{Orbit Equivalence for Cantor Minimal $\Z^d$-systems}, Invent. Math. 179 (2010), no. 1, 119--158.
\bibitem[GPS95]{Giordano1995} Thierry Giordano, Ian F. Putnam, and Christian F. Skau, \emph{Topological orbit equivalence and $C^*$-crossed products}, J. Reine Angew. Math. 469 (1995), 51--111.
\bibitem[GPS04]{Giordano2004} Thierry Giordano, Ian F. Putnam, and Christian F. Skau, \emph{Affable equivalence relations and orbit structure of Cantor dynamical systems}, Ergodic Theory Dyn. Syst. 24 (2004), no. 2, 441--475.
\bibitem[GW95]{Glasner1995a} Eli Glasner and Benjamin Weiss, \emph{Weak orbit equivalence of Cantor minimal systems}, Int. J. Math. 6 (1995), no. 4, 559--579.
\bibitem[HKY11]{Keane2011} Toshihiro Hamachi, Michael S. Keane, and Hisatoshi Yuasa, \emph{Universally measure-preserving homeomorphisms of Cantor minimal systems}, J. Anal. Math. 113 (2011), 1--51.
\bibitem[Kri80]{Krieger1979} Wolfgang Krieger, \emph{On a dimension for a class of homeomorphism groups}, Math. Ann. 252 (1980), no. 2, 87--95.
\bibitem[Mat08]{Matui2008} Hiroki Matui, \emph{An absorption theorem for minimal AF equivalence relations on Cantor sets}, J. Math. Soc. Japan 60 (2008), no. 4, 1171--1185.
\bibitem[MR23]{Melleray2023} Julien Melleray and Simon Robert, \emph{From invariant measures to orbit equivalence, via locally finite groups}, Ann. Henri Lebesgue 6 (2023), 259--295.
\bibitem[Put10]{Putnam2010} Ian F. Putnam, \emph{Orbit equivalence of Cantor minimal systems: a survey and a new proof}, Expo. Math. 28 (2010), no. 2, 101--131.
\bibitem[Put18]{Putnam2018} Ian F. Putnam, \emph{Cantor minimal systems}, University Lecture Series, vol. 70, American Mathematical Society, 2018.
\end{thebibliography}
\end{document}
